\section{Open Questions}

The following five open research questions arose specifically from doing
\emph{this} replication (not generic future-work items copied from the
paper).

\begin{description}
\item[Q1 --- Detector-confusion sensitivity of the Bell-basis advantage.]
Under a realistic per-detector confusion / miscalibration model on the four
Bell-state detectors $D_1..D_4$ (i.e., off-diagonal elements in the POVM),
does the Bell-basis VQE's advantage in shot budget survive, or does the
coupled error on $\{\braket{XX},\braket{YY},\braket{ZZ}\}$ amplify the
energy variance enough to erase it? The paper's Fig.~3 asserts partial
error mitigation for certain Hamiltonians but does not test this
systematically across Hamiltonians; our replication did not implement a
detector-noise model.

\emph{Basis:} the main VQE-E result reproduced exactly under multinomial
(ideal) Bell sampling. We never added a noisy-detector POVM, so the
paper's error-mitigation claim (C6) remains untested in this
replication.

\item[Q2 --- HeH$^+$ Hamiltonian factor-of-2 unit convention.]
The paper's Appendix-A HeH$^+$ coefficient table (Fig.~A1) and its
``$E_\mathrm{th}=-2.863$~MJ/mol'' at $R=0.9$~\AA{} differ by a
\emph{clean} factor of exactly 2 when compared to the ground eigenvalue of
$H=\sum_j w_j\,\sigma_j$ built from those tabulated coefficients (we found
$E_\mathrm{exact} = -5.7252$~MJ/mol). Which is the intended normalization
--- is the paper's table pre-divided by 2, is $E_\mathrm{th}$ reported
per electron pair, or is one of the two just a unit ambiguity between
MJ/mol and Ha$\times N_A$?

\emph{Basis:} Building $H$ directly from Fig.~A1 row $R=0.9$~\AA{} gives
$-5.7252$~MJ/mol; the paper says $E_\mathrm{th}=-2.863$~MJ/mol. Ratio =
1.9997 (within 0.02\%). Our best VQE ($-5.7259$~MJ/mol using the paper's
grouping) matches our exact reference to $<0.001$~MJ/mol.

\item[Q3 --- Scaling of GC grouping advantage.]
For H$_2$/STO-3G the greedy general-commutativity grouping collapsed 15
Pauli strings into just 2 groups (all-Z-diagonal + all XXYY-parity). But
this collapse is possible only because 4-qubit fermionic parity gives
especially large commuting families; how does this basis-count scaling
behave as we go to H$_2$ in a bigger basis (6-31G, 4-orbital, 8-qubit) or
LiH (12-qubit)? Does the Bell/GC advantage stay $>30\%$ or shrink, and is
greedy-GC still near-optimal or does one need the harder graph-coloring
formulation?

\emph{Basis:} H$_2$/STO-3G at 4 qubits was a `lucky' Hamiltonian: only two
commuting families sufficed. The paper's cases are only 2-qubit; scaling
behaviour is not addressed.

\item[Q4 --- COBYLA-tolerance sensitivity of Bell VQE.]
In the Heisenberg VQE-E run under loose COBYLA \texttt{tol=0.01} we
observed much larger cross-run variance ($\sigma\!\sim\!0.78$) than VQE-P
($\sigma\!\sim\!0.02$), even though each Bell measurement returned a
lower-variance $\braket{H}$ estimate at fixed shots (correlated estimates
of $\braket{XX},\braket{YY},\braket{ZZ}$ from one basis). Tightening to
$\texttt{tol}=0.001$ restored VQE-E to VQE-P quality. Is the
COBYLA-tolerance $\to$ optimizer-early-stopping interaction with
correlated shot noise a systematic warning for practitioners who read the
paper's \texttt{tol=0.01} as canonical, or is it specific to our ansatz
choice?

\emph{Basis:} Directly observed in R2; same shot budget, same ansatz, only
difference is optimizer tolerance. VQE-E was harmed more by loose tol
than VQE-P.

\item[Q5 --- Per-iteration vs total-shot budget in the paper.]
The paper reports ``total number of shots is fixed to be $N=9000$'' without
saying how many shots the classical optimizer used per iteration vs how
many trajectories were used to compute the reported std bar. Our VQE-P and
VQE-E converged with $\sim50$--$75$ COBYLA function evaluations at 3000
(resp.\ 9000) shots per evaluation; if the paper's $N=9000$ is total
across all iterations rather than per iteration, then their per-eval
shots must be $\ll 100$ and the energy estimator would be shot-noise
dominated. Which interpretation is intended?

\emph{Basis:} We tried to match the paper's shot protocol and had to make
an interpretation choice. Both interpretations reproduce the qualitative
result, but they imply very different practical shot budgets in an
experiment.
\end{description}
